\clearpage
\VerdierPhysicalPageAnchor{0240}%
\begingroup\small\setlength{\parskip}{0pt}\setlength{\emergencystretch}{1.5em}
\setlength{\abovedisplayskip}{4pt plus 2pt minus 2pt}%
\setlength{\belowdisplayskip}{4pt plus 2pt minus 2pt}%

\VerdierSemanticAnchor{ch4d-4-4-5}
\noindent\textbf{4.4.5.}\quad
Les données sont les mêmes qu'en (4.4.4) mais on se donne, de plus, un
objet $X$ de $D(\mathcal A)$. Soit
$\bigl(X(p,q),\delta(p,q,r)\bigr)$ le deuxième objet spectral à valeurs
dans $D(\mathcal A)$ associé au complexe $X$ (4.3.1). Le foncteur $F$
transforme cet objet spectral en un objet spectral à valeurs dans
$\mathcal B$. On en déduit une suite spectrale fonctorielle en $F$ et en
l'argument $X$ lorsque celui-ci varie dans la catégorie $D(\mathcal A)$
(4.3.3). Cette suite spectrale est appelée la \emph{deuxième suite
spectrale du foncteur $F$, relative au complexe $X$}. On utilise toujours,
pour cette suite spectrale, l'indexation introduite au (chap.~II, 4.3.3.2).

\VerdierSemanticAnchor{ch4d-4-4-6}
\medskip\noindent\textbf{Proposition 4.4.6.}\quad

\VerdierSemanticAnchor{ch4d-4-4-6-a}
\noindent\textit{a) Le terme $\mathrm{II}_2^{p,q}$ de la deuxième suite
spectrale de $F$ relative au complexe $X$ est donné par :}
\[
\mathrm{II}_2^{p,q}=F^p\bigl(H^q(X)\bigr).
\]

\VerdierSemanticAnchor{ch4d-4-4-6-b}
\noindent\textit{b) Les termes $\mathrm{II}_{\infty}^{p,q}$ sont les
gradués associés aux objets $F^{p+q}(X)$ convenablement filtrés.}

\VerdierSemanticAnchor{ch4d-4-4-6-proof}
En effet (4.3.4.1), on a un isomorphisme canonique :
\[
X(p,p+1)\xrightarrow{\sim}H^p(X)[-p],
\]
d'où en utilisant (chap.~II, 4.3.3.2) :
\[
\mathrm{II}_2^{p,q}
=F^{p+q}\bigl(H^q(X)[-q]\bigr)
=F^p\bigl(H^q(X)\bigr),
\]
d'où l'égalité (abusive) de a). L'assertion b) résulte de
(chap.~II, 4.3.2).

\VerdierSemanticAnchor{ch4d-4-5}
\subsection{Problèmes de convergence.}

\VerdierSemanticAnchor{ch4d-4-5-1}
\noindent\textbf{Définition 4.5.1.}\quad
Soient $\mathcal A$ et $\mathcal B$ deux catégories abéliennes. Un foncteur
cohomologique $F\colon D^*(\mathcal A)\longrightarrow\mathcal B$
($*=+,-,b,{}$ ``vide'') est dit \emph{stationnaire à droite} (resp.\
\emph{à gauche}) s'il existe un entier $n_0$ tel que, pour tout objet $X$
de $D^*(\mathcal A)$ tel que $H^i(X)=0$, pour tout $i<0$ (resp. $>0$), on
ait $F^n(X)=0$, pour tout $n<n_0$ (resp. $n>n_0$). Le foncteur $F$ est dit
\emph{stationnaire} s'il est stationnaire à droite et à gauche. Un foncteur
exact $G\colon D^*(\mathcal A)\longrightarrow D^{*'}(\mathcal B)$ est dit
\emph{stationnaire à droite} (resp.\ \emph{stationnaire à gauche}, resp.\
\emph{stationnaire}) si le foncteur cohomologique obtenu en composant le
foncteur $G$ avec le foncteur cohomologique canonique
$D^{*'}(\mathcal B)\longrightarrow\mathcal B$ est stationnaire à droite
(resp.\ stationnaire à gauche, resp.\ stationnaire).
\endgroup
